\documentclass[11pt]{article}

\usepackage[T1]{fontenc}
\usepackage[utf8]{inputenc}
\usepackage{lmodern}
\usepackage[margin=1in]{geometry}
\usepackage{microtype}
\usepackage{amsmath}
\usepackage{booktabs}
\usepackage{tabularx}
\usepackage{threeparttable}
\usepackage{multirow}
\usepackage{makecell}
\usepackage{placeins}
\usepackage{enumitem}
\usepackage{caption}
\usepackage{xcolor}
\usepackage{pgfplots}
\usepgfplotslibrary{groupplots}
\pgfplotsset{compat=1.18}
\pgfplotstableset{col sep=semicolon}
\usepackage[numbers,sort&compress]{natbib}
\usepackage[hidelinks]{hyperref}
\usepackage[type={CC},modifier={by},version={4.0}]{doclicense}

\setlist{itemsep=2pt,topsep=4pt}
\captionsetup{font=small,labelfont=bf}
\newcommand{\arm}[1]{\texttt{#1}}
\newcolumntype{L}{>{\raggedright\arraybackslash}X}
\newcommand{\pholm}{p_{\mathrm{Holm}}}
\newcommand{\STUDYFOURRUNS}{330}
\newcommand{\STUDYFOURCOST}{75.07}

\title{What's in a Name?\\[4pt]\large Addressability, Parental Framing and Group Labels in Generational Populations of LLM Agents}
\newcommand{\PAPERDOI}{10.5281/zenodo.23397390}
\newcommand{\SOFTWAREDOI}{10.5281/zenodo.23397394}
\author{Vaughan Wynne-Jones\\{\small Independent researcher}\\
{\small ORCID: \href{https://orcid.org/0009-0006-3469-7151}{0009-0006-3469-7151}}}
\date{October 2026\ifx\PAPERDOI\empty\else\\[2pt]{\small Preprint, \href{https://doi.org/\PAPERDOI}{doi:\PAPERDOI}}\fi}

\begin{document}
\maketitle

\begin{abstract}
Names in multi-agent simulations of language-model agents are usually treated as scaffolding.
We varied them as the experimental manipulation in a generational population of LLM agents
(Claude Haiku~5.5). The agents forage, hunt in pairs, form couples, endow their children, compete
with them for survival, and vote on their own scoring rule. Every arm saw identical information
and identical random draws. Only the handle each agent carried changed, and how that handle was
described. The evidence comes from four pre-registered studies: 480 analysed eight-generation
runs and about 367,000 model calls.

Any stable individual handle, whether a name or a bare ID and whoever gave it, raised pair
hunting two- to fivefold over no handle. Couple formation instead followed a parental frame. A
meaningless ID raised it when parents picked it, and equally when agents were merely told their
parents had given it. Only names composed by the parents raised investment in newborns, and
hunters with such names shared their kills most often. Arbitrary Blue/Green bands produced
in-group pairing and hunting. Yet a band added to an individual ID doubled pair hunting rather
than splitting agents into factions.

Set against human findings registered in advance, the identification and group-label effects
matched in direction. Two did not match. Namesaking did not raise parental investment, and agents
put about 3\% of their holdings into a public good, against the 40--60\% typical of people.

Exploratory analyses add three leads:
\begin{itemize}
  \item a world-assigned ID lowered couple formation below no handle at all;
  \item parents gave less to a child that the lot placed in the other band;
  \item a single sentence saying that children are named after a parent halved couple formation.
\end{itemize}

How agents are named, and how naming is described, shapes cooperation and family formation in
LLM agent populations.
\end{abstract}


\section{Introduction}

Populations of language-model agents are now used to study cooperation, norms and governance
\citep{park2023generative,vallinder2024cultural,akkil2026emergence,akkil2026stress}. In most of
these systems an agent's name is scaffolding: a label the designer attaches so that agents and
logs can refer to one another. Human names do more than that. A name lets others point at the
same individual again, which is what reputation and repeated cooperation need. A name is also
given by someone, usually a parent, and it can carry family meaning: a child named after a parent
advertises the tie \citep{obasi2016naming}, and in at least one human sample such children receive
more paternal investment \citep{saffa2021namesake}.

This paper asks whether those two functions can be separated in a population of LLM agents, and
which behaviours follow which. We built a generational simulation in which agents forage, hunt in
pairs, guard one another's stores, form couples, have children, decide how much to give each
newborn, compete with their own children for a place in the next generation, and vote on the
scoring rule. Every agent is the same model, Claude Haiku~5.5. The only thing that differs
between the arms of the experiment is the handle each agent carries and how it is described: a
name chosen by the agent's parents, a name or an ID assigned by the world, an ID picked by
parents, a group label, or nothing at all. Everything else, including the history every agent
sees of its dealings with every other agent, is identical across arms and driven by common random
numbers.

We report four pre-registered studies run in sequence, each designed in response to the last:

\begin{enumerate}
  \item \textbf{Named vs unnamed.} Parent-named populations formed more couples and held about
        five times as many pair hunts as populations without names.
  \item \textbf{A stable ID.} A world-assigned ID restored pair hunting to the named level but
        did nothing for couple formation, which separated \emph{addressability} from
        \emph{naming}.
  \item \textbf{A $2\times2$ of who gives the handle by what it is.} Having parents give the
        handle raised couple formation even when the handle was a meaningless ID; a word name
        raised it only when the world assigned it; and only names chosen by parents raised
        investment in newborns.
  \item \textbf{A novelty round} on fresh seeds with eleven arms. The earlier results replicated,
        and kill sharing followed parental names. Either the parental act or merely its wording
        raised couple formation. Naming a child after a parent did not raise that parent's
        investment. Arbitrary group labels biased pairing and hunting toward the in-group but,
        added to individual IDs, doubled coordination. Agents put almost nothing into a shared
        pot.
\end{enumerate}

Each test in the fourth study was registered with the human finding it is meant to be compared
with, so that the results can be set beside the human literature (Section~\ref{sec:human}).
The comparisons are of direction, not magnitude: these are claims about one model under one
framing, not about people.

\section{Related work}
\label{sec:related}

\paragraph{Generational and social LLM agent populations.}
Generative agents showed that LLM-driven characters can sustain plausible social life in a shared
world \citep{park2023generative}. Long-horizon agent worlds in which agents govern themselves
through democratic mechanisms, the starting point for this project, are studied on the Emergence
World platform \citep{akkil2026emergence,akkil2026stress}. \citet{vallinder2024cultural} evolved
generational populations playing the Donor Game and found that societies of Claude~3.5 Sonnet
agents sustained markedly more cooperation than those of Gemini~1.5 Flash or GPT-4o. Two recent simulations add reproduction: \citet{zhou2026moral} let agents
with prompted moral types reproduce, and found that kin-focused agents prioritise offspring;
\citet{pasichnyk2026yerkes} builds reproduction on parental-investment theory, with offspring
inheriting mutated attributes. Neither varies names.

\paragraph{Identity labels in multi-agent LLM systems.}
\citet{delgiudice2026identity} found that showing agents a model-family label, or even an
arbitrary shared label, split them into factions and degraded coordination; withholding the
labels removed the effect. \citet{chen2026favor} found that a group label alone produces little
in-group favouritism when giving is costly, and that shared interaction history matters more.
\citet{hu2026dissociative} argue that reputation mechanisms presuppose persistent identities,
which LLM agents often lack. Those studies vary \emph{group} labels or the persistence of
identity across sessions. We vary unique \emph{individual} handles, who gives them and what they
say, inside one persistent world, and in the fourth study set group labels beside them directly.

\paragraph{Human research.}
In public-goods experiments, identifying individuals raises contributions, while identifying only
groups does little \citep{andreoni2004public,christens2019identification}; names alone had no
direct effect in one recent study \citep{sampaio2023identity}. Standard public-goods games find
contributions starting around 40--60\% of the endowment and declining with repetition
\citep{ledyard1995public,chaudhuri2011sustaining,fehr2000cooperation}. Arbitrary group labels are
enough to produce in-group favouritism \citep{tajfel1971social}. Kin selection and
parental-investment theory predict investment that tracks relatedness and certainty of parentage
\citep{hamilton1964genetical,trivers1972parental}, and namesaking has been studied as a signal of
that tie \citep{obasi2016naming,saffa2021namesake}.

As far as a literature search found, no earlier study varies personal names in a generational,
reproducing population of LLM agents while holding everything else fixed.

\section{The simulation}
\label{sec:sim}

The simulator is open source \citep{repo2026agentnames}. Its core holds the rules and knows
nothing about the model, storage or random source; each arm is a configuration of the same core.

\subsection{World and generation loop}

A run starts with 10 founders, each a pure specialist in one of three resources: energy, food and
memory. Every agent has three skills from 0 to 1, one per resource. A run lasts 8 generations,
and each generation proceeds as follows.

\begin{enumerate}
  \item \textbf{Archivist offer.} Every agent with ancestors or descendants on record is offered an
        account of them from the archivist, at a cost of 10\% of its first-round yield.
  \item \textbf{Foraging.} Five rounds. In each round every agent chooses one activity
        (Table~\ref{tab:activities}), may give resources to others, and may post one public
        message that everyone sees at their next decision.
  \item \textbf{The commons} (Study~4 only). Every agent may put resources into a shared pot,
        which is multiplied by 4 and split equally among the 10 agents. Everyone then sees what
        each agent put in.
  \item \textbf{Pairing.} Each agent ranks up to three others. A couple forms only when both chose
        each other, best combined rank first, up to five couples. Close kin cannot pair.
  \item \textbf{Birth.} Each couple has one child, who inherits each skill whole from one parent
        or the other, with a 5\% chance per skill of mutation. Each parent then decides, without
        seeing the co-parent's choice, how much of its own holdings to give the child, which
        starts with nothing else.
  \item \textbf{Scoring and the cut.} Everyone present, newborns included (up to 15), is scored
        on what they hold. Only the top 10 stay. Parents and children compete in the same
        scoring, so a parent can be displaced by its own child.
  \item \textbf{Governance.} Each survivor may propose one change to the scoring rule: the weights
        on each resource, a \emph{kin weight} that adds relatives' scores weighted by
        relatedness, a \emph{group weight} that blends an agent's score with the mean, and a
        credit for net giving. Up to five proposals go to a vote; a majority of survivors passes
        one.
\end{enumerate}

Resources reset each generation. Each generation's highest-scoring agent, and the final
survivors, enter a permanent archive, which every agent is told about from the start.

\begin{table}[t]
\centering
\caption{Foraging activities. Every yield varies by $\pm20\%$, and even a skill of 0 gives a fifth
of the full-skill yield.}
\label{tab:activities}
\begin{tabularx}{\linewidth}{@{}lX@{}}
\toprule
Activity & What it does at full skill \\
\midrule
\arm{gather\_energy} & 3 energy, alone. \\
\arm{hunt\_alone} & 1.5 food, alone. \\
\arm{hunt\_together} & 10 food in total, with an $80\%\times$skill chance of success, but only if
both partners name each other in the same round. One hunter, drawn by skill, lands the kill and
holds all of it; whether it is shared is up to that hunter. \\
\arm{collect\_memory} & 3.5 memory. Collecting does not protect the store. \\
\arm{guard\_memory} & Nothing for the guard; protects another agent's store. An unguarded store
loses 30\% of its start-of-round contents each round. \\
\bottomrule
\end{tabularx}
\end{table}

\subsection{What agents see}

Each decision is a separate model call. The system prompt describes the world and is identical
for every call in a run. The user message gives the agent's own skills, holdings and standing;
what happened since its last decision; the last round's public messages; and a line for every
other agent present. That line gives the other agent's relation to the decider (``your elder
parent'', ``your cousin'', ``no known relation''), its skills and holdings, and the decider's
history with it: what each gave the other, guarding in each direction, and hunts together.

This history is given in every arm. The experiment therefore tests whether handles change
behaviour when information about others is equal, not whether handles make tracking possible.
In every arm other agents are listed as \texttt{[1]}\ldots\texttt{[n]}, and the numbers are
reshuffled at every prompt so that they cannot harden into names. The handle, where there is
one, is shown after the number.

\subsection{Arms}

Table~\ref{tab:arms} lists the eleven arms. The arm changes only the handle, how it is described,
how the archivist phrases its answers, and what the archive keeps. Founders in the word-name arms
carry fixed names (Ash, Brook, Cedar, \ldots); in the ID arms, IDs have the form \texttt{A47},
drawn from a seed-shuffled pool so that numbers do not reveal age. ID arms are told that an ID is
kept for life and never shared; name arms are not told this about names.

\begin{table}[t]
\centering
\small
\caption{Arms. Studies 1--3 used the first five; Study~4 used all eleven.}
\label{tab:arms}
\begin{tabularx}{\linewidth}{@{}l X c@{}}
\toprule
Arm & Handle, and what parents do & Study \\
\midrule
\arm{named} & A word name. Both parents propose names for each child; a lot decides whose list is used. & 1--4 \\
\arm{unnamed} & No handle. & 1--4 \\
\arm{stable\_id} & An ID assigned by the world. & 2--4 \\
\arm{world\_name} & A single-word name assigned by the world from a fixed pool. & 3--4 \\
\arm{parent\_id} & An ID picked by the parents from three the world offers; agents are told parents gave it. & 3--4 \\
\midrule
\arm{parent\_id\_plain} & As \arm{parent\_id}, but no prompt says parents gave the ID. & 4 \\
\arm{stable\_id\_parental} & As \arm{stable\_id}, but agents are told their parents gave them their IDs. & 4 \\
\arm{world\_name\_shown} & As \arm{world\_name}, with the child's name shown to parents when they decide the endowment. & 4 \\
\arm{world\_namesake} & The world names each child after one parent, chosen by lot (Fern~II), and tells each parent ``after you'' or ``after your co-parent''. & 4 \\
\arm{band\_only} & An arbitrary Blue or Green band, assigned by lot, and no individual handle. & 4 \\
\arm{stable\_id\_banded} & A world-assigned ID plus a band. & 4 \\
\bottomrule
\end{tabularx}
\end{table}

\subsection{Common random numbers}

Every random draw is a hash of the run's seed and a key that names the draw, such as the
generation, birth order and skill slot. Draws never depend on call order. Two arms run on the same
seed therefore get the same inheritance and mutation outcomes for the same birth, even when their
histories differ, and a scripted agent that ignores handles plays out an identical world in every
arm. Both properties are covered by the test suite, along with checks that each arm's prompts
contain no handle of another arm's kind.

\section{Methods}
\label{sec:methods}

\paragraph{Model.} Every agent in every arm is Claude Haiku~5.5 (\texttt{claude-haiku-5-5}) at
effort \texttt{low}, called through the Anthropic API with structured JSON output. The system
prompt is cached. A decision that fails to parse after a retry is replaced by a logged default
(a \emph{fallback}). A run of 8 generations makes about 700--800 calls and costs about \$0.20--0.25.

\paragraph{Measures.} The registered measures are:
\begin{itemize}
  \item \textbf{pair hunts}: pair hunts that took place in a run;
  \item \textbf{couples per generation}: couples formed, averaged over generations;
  \item \textbf{endowment fraction}: the share of its holdings a parent gives its newborn, averaged
        over parents' birth decisions;
  \item \textbf{gathered per forager}: resources gathered per agent per generation;
  \item \textbf{shared kills}: successful pair hunts after which the hunter who landed the kill
        gave its partner food later in the same generation;
  \item \textbf{commons contribution share}: the share of its holdings an agent puts into the pot;
  \item \textbf{in-band bias}: for gifts, guards, hunts or couples, the share directed within the
        actor's own band, minus the share of same-band agents among the actor's living others at
        the time (zero under no bias);
  \item \textbf{namesake differences}: in \arm{world\_namesake}, the endowment share, and later
        gifts, from the parent a child is named after minus those from its co-parent.
\end{itemize}

\paragraph{Tests.} Every registered test is a two-sided sign-flip test over seeds: paired by seed
for arm-versus-arm contrasts, against zero for one-sample tests. Up to 20 seeds the test is exact;
beyond that it uses 100,000 random sign patterns from a fixed seed, with $p=(k+1)/(N+1)$, so the
smallest reportable $p$ is about $10^{-5}$. Holm--Bonferroni correction is applied within each
registered family at $\alpha=0.05$. A hypothesis counts as supported when its Holm-adjusted $p$ is
below 0.05 and the effect goes the predicted way.

\paragraph{Exclusions.} A run is excluded if it fails to complete or if more than 5\% of its
decisions are fallbacks. In Study~1 an excluded run's seed partner was dropped and the pair
replaced by a fresh seed; from Study~2 on, an excluded seed is re-run once and, if it fails again,
dropped from every test that uses that arm.

\paragraph{Pre-registration.} Each study's hypotheses, tests, correction, exclusion rule and
configuration were committed to the public repository before any run on its seeds
(Table~\ref{tab:studies}). Analyses outside the registered tests are labelled exploratory and
reported unadjusted. Three development pilots (32 runs in all) preceded Study~1 and shaped the
mechanics; none of their data enter the registered tests.

\begin{table}[t]
\centering
\small
\caption{The four studies. Studies 2 and 3 reuse Study~1's seeds and are paired with its runs by
seed; Study~4 uses fresh seeds and runs all arms in one batch.}
\label{tab:studies}
\begin{tabular}{@{}l l r r r l@{}}
\toprule
Study & New arms & Seeds & Runs & Cost (USD) & Registered at \\
\midrule
1 & \arm{named}, \arm{unnamed} & 30 & 64$^a$ & 12.07 & \texttt{a324dbc} \\
2 & \arm{stable\_id} & 30 & 30 & 5.88 & \texttt{6136b3e} \\
3 & \arm{world\_name}, \arm{parent\_id} & 30 & 60 & 12.19 & \texttt{0b92bed} \\
4 & all eleven arms, plus the commons & 30 & \STUDYFOURRUNS & \STUDYFOURCOST & \texttt{bc83a2e} \\
\bottomrule
\end{tabular}
\par\smallskip
{\footnotesize $^a$ Four runs stopped at the account's API usage limit; both affected seed pairs
were replaced as registered.}
\end{table}

\section{Results}
\label{sec:results}

\subsection{Study 1: names raise coordination}

Across 30 seed pairs, named populations formed more couples per generation than unnamed ones
(2.38 vs 1.74; named higher in 23 of 30 pairs; $\pholm=0.00008$) and held about five times as
many pair hunts (13.4 vs 2.7; 28 of 30; $\pholm=0.00003$). The third registered measure,
gathering per forager, leaned the same way but was not significant (14.01 vs 13.80; $\pholm=0.11$).

Exploratory measures suggested that named parents invested more in their children: they gave a
larger share of their holdings at birth (0.55 vs 0.47, $p=0.003$), more of their newborns
survived the first cut, and more of them were displaced by their own child. The design doc's
tribalism prediction, that names plus visible lineage would concentrate help inside families,
was not borne out: unnamed agents gave 84\% of their gift amounts to kin, named agents 44\%.

A name, however, is also a stable handle. Named agents can say ``Brook, let's hunt this round''
and mean the same agent next round; unnamed agents can only use bracket numbers that are
reshuffled before the next decision. Study~1 cannot say whether the effects come from being named
by a parent or from having any stable handle. (On Study~4's fresh seeds the couples difference
between named and unnamed populations was smaller and not significant; see
Section~\ref{sec:study4}.)

\subsection{Study 2: addressability and naming come apart}

A world-assigned stable ID restored pair hunting to the named level (stable ID vs unnamed: 13.1
vs 2.7; 26 of 30; $\pholm=0.00004$; named vs stable ID: 13.4 vs 13.1; $\pholm=0.87$). It did
nothing for couples: named populations formed far more couples than stable-ID ones (2.38 vs 1.44;
27 of 30; $\pholm=0.00004$), and stable IDs gave no lift over no handle at all (1.44 vs 1.74;
$\pholm=0.092$). Being able to point at the same agent again is what lets agents coordinate a
hunt. Choosing partners and investing in children depend on something a parent-given name has and
a bare ID lacks.

\subsection{Study 3: who gives the handle, and what it is}

Study~3 completed a $2\times2$ of who chooses the handle (parents or world) by what it is (a word
or an ID). Eight contrasts were registered with one Holm family (Table~\ref{tab:twobytwo}).

\begin{table}[t]
\centering
\small
\caption{Study 3: registered contrasts, paired over 30 seeds, Holm-corrected across all eight.}
\label{tab:twobytwo}
\begin{tabular}{@{}l l c c c c@{}}
\toprule
Contrast & Measure & Means & First higher & $\pholm$ & Supported \\
\midrule
Parent ID vs stable ID & Couples & 2.53 vs 1.44 & 29/30 & 0.00008 & yes \\
World name vs stable ID & Couples & 2.01 vs 1.44 & 23/30 & 0.0005 & yes \\
Named vs parent ID & Endowment & 0.55 vs 0.47 & 19/30 & 0.037 & yes \\
Named vs world name & Couples & 2.38 vs 2.01 & 21/30 & 0.039 & yes \\
Named vs world name & Endowment & 0.55 vs 0.46 & 19/30 & 0.058 & no \\
Named vs parent ID & Couples & 2.38 vs 2.53 & 11/30 & 0.50 & no \\
Parent ID vs stable ID & Endowment & 0.47 vs 0.43 & 19/30 & 0.50 & no \\
World name vs stable ID & Endowment & 0.46 vs 0.43 & 17/30 & 0.50 & no \\
\bottomrule
\end{tabular}
\end{table}

Three things stand out (Figure~\ref{fig:early}).
\begin{itemize}
  \item \textbf{Parents giving the handle raised couple formation, whatever the handle.} The
        largest effect in the study was parent-picked IDs against world-assigned IDs. The
        ``choice'' was nominal: parents took the first of the three offered IDs 96\% of the time,
        and both parents picked the same one in 94\% of births. The lift comes from the parental
        role, which agents are told about from the start and enact at every birth, not from
        deliberating over a label.
  \item \textbf{A word name helped only when the world assigned it.} Among world-assigned handles,
        names formed more couples than IDs; among parent-given handles, names added nothing.
  \item \textbf{Investment in newborns rose only with names chosen by parents.} Named-arm parents
        gave about 0.55 of their holdings; every other arm gave 0.43--0.47, including unnamed.
\end{itemize}

Pair hunting, an exploratory measure here, was high in all four handle arms (10.8--13.4) and low
only without a handle (2.7), confirming Study~2.

\begin{figure}[t]
\centering
\begin{tikzpicture}
\begin{groupplot}[
  group style={group size=3 by 1, horizontal sep=1.6cm},
  width=0.345\linewidth, height=5.2cm,
  ybar, /pgf/bar width=9pt, ymin=0,
  xtick=data, xticklabels from table={figdata/early.csv}{label},
  xticklabel style={rotate=45, anchor=east, font=\footnotesize},
  yticklabel style={font=\footnotesize},
  title style={font=\small},
  enlarge x limits=0.15,
  error bars/y dir=both, error bars/y explicit,
  cycle list={{fill=black!30, draw=black!80}},
]
\nextgroupplot[title={Pair hunts per run}]
\addplot table[x=x, y=pair_hunts, y error=pair_hunts_se] {figdata/early.csv};
\nextgroupplot[title={Couples per generation}]
\addplot table[x=x, y=couples, y error=couples_se] {figdata/early.csv};
\nextgroupplot[title={Share given to a newborn}, ymax=0.65]
\addplot table[x=x, y=endowment, y error=endowment_se] {figdata/early.csv};
\end{groupplot}
\end{tikzpicture}
\caption{Studies 1--3, five arms on the same 30 seeds: means over seeds with standard errors.
Pair hunting follows any stable handle; couple formation follows a parent-given handle or a word
name; investment in newborns stands out only with parent-chosen names.}
\label{fig:early}
\end{figure}

\paragraph{Leads carried into Study 4.} After Study~3 we tabulated every action by arm
(exploratory). Two patterns were new. First, hunters in the word-name arms shared their kills
more often: the hunter who landed a pair-hunt kill later gave its partner food after about 21\%
of successful hunts in \arm{named} and 20\% in \arm{world\_name}, against 10\% in
\arm{parent\_id}, 6\% in \arm{stable\_id} and 5\% in \arm{unnamed}. Second, unnamed agents,
when they gave at all, gave mostly to kin. Study~4 registers the first as a test.

\FloatBarrier
\subsection{Study 4: the novelty round}
\label{sec:study4}

Study~4 ran all eleven arms on 30 fresh seeds (201--230) in one batch, so every comparison is
within one run window. It added the commons to every arm. All 330 runs completed without a single
fallback decision in 258,808 calls, so nothing was excluded. Nineteen tests were registered in
five families, each corrected separately (Table~\ref{tab:round5}); Figure~\ref{fig:round5} shows
the main measures in every arm.

\begin{table}[tbp]
\centering
\footnotesize
\caption{Study 4: registered tests over 30 fresh seeds, Holm-corrected within each family. For
arm-versus-arm tests, ``Result'' gives the two means and the number of seeds on which the first
arm was higher; for one-sample tests, the mean over seeds and the number above zero. Tests marked
``different'' were registered without a predicted direction.}
\label{tab:round5}
\begin{tabularx}{\linewidth}{@{}l L l l l l@{}}
\toprule
& Comparison and measure & Result & $p$ & $\pholm$ & Verdict \\
\midrule
\multicolumn{6}{@{}l}{\textit{A. Replication and kill sharing}} \\
A1 & Named vs unnamed: pair hunts & 12.7 vs 3.5; 26/30 & 0.00003 & 0.0002 & supported \\
A2 & Stable ID vs unnamed: pair hunts & 8.4 vs 3.5; 23/30 & 0.0009 & 0.004 & supported \\
A3 & Parent ID vs stable ID: couples & 2.22 vs 1.40; 23/30 & 0.00004 & 0.0002 & supported \\
A4 & Named vs parent ID: endowment & 0.50 vs 0.41; 22/30 & 0.005 & 0.014 & supported \\
A5 & Named vs parent ID: shared kills & 1.93 vs 0.97; 15/30$^a$ & 0.010 & 0.020 & supported \\
A6 & World name vs stable ID: shared kills & 0.97 vs 0.37; 13/30$^a$ & 0.0498 & 0.0498 & supported$^b$ \\
\midrule
\multicolumn{6}{@{}l}{\textit{B. The parental act or its wording}} \\
B1 & Parent ID, plain wording vs stable ID: couples & 2.02 vs 1.40; 22/30 & 0.0006 & 0.0006 & supported \\
B2 & Stable ID, parental wording vs stable ID: couples & 2.08 vs 1.40; 22/30 & 0.0002 & 0.0005 & supported \\
\midrule
\multicolumn{6}{@{}l}{\textit{C. Namesaking}} \\
C1 & Namesake parent minus co-parent: endowment & $-0.015$; 15/30 & 0.54 & 0.54 & not supported \\
C2 & Namesake vs world name shown: endowment & 0.43 vs 0.34; 21/30 & 0.021 & 0.064 & not supported \\
C3 & Namesake parent minus co-parent: later gifts & $+0.73$; 8 above, 3 below & 0.035 & 0.069 & not supported \\
\midrule
\multicolumn{6}{@{}l}{\textit{D. Group labels and individual handles}} \\
D1 & Both band arms: in-band couple bias & $+0.13$; 26/30 & 0.00002 & 0.0001 & supported \\
D2 & Both band arms: in-band gift bias & $+0.16$; 5/6$^c$ & 0.094 & 0.19 & not supported \\
D3 & ID with band: in-band hunt bias & $+0.20$; 21/30 & 0.005 & 0.016 & supported \\
D4 & ID with band vs stable ID: pair hunts & 16.7 vs 8.4; 22/30 & 0.0006 & 0.002 & different \\
D5 & Band only vs unnamed: pair hunts & 3.1 vs 3.5; 11/30 & 0.71 & 0.71 & no difference \\
\midrule
\multicolumn{6}{@{}l}{\textit{E. The commons}} \\
E1 & Stable ID vs unnamed: contribution share & 0.033 vs 0.028; 11/30 & 0.66 & 1.0 & not supported \\
E2 & Band only vs unnamed: contribution share & 0.029 vs 0.028; 14/30 & 0.90 & 1.0 & no difference \\
E3 & Named vs stable ID: contribution share & 0.043 vs 0.033; 15/30 & 0.53 & 1.0 & no difference \\
\bottomrule
\end{tabularx}
\par\smallskip
{\footnotesize $^a$ Nine (A5) and twelve (A6) seeds had no shared kill in either arm.
$^b$ The margin below 0.05 is smaller than the Monte Carlo error of $p$ (about $\pm0.0007$).
$^c$ Only six seeds had gifts in both band arms.}
\end{table}

\paragraph{The earlier results replicate.} On fresh seeds, any stable handle again raised pair
hunting (A1, A2), parent-picked IDs again formed more couples than world-assigned IDs (A3), and
named-arm parents again gave the largest share of their holdings to newborns (A4). Kill sharing,
an exploratory lead from Studies 1--3, was confirmed for both registered contrasts, although A6
only narrowly. Per successful pair hunt, the hunter who landed the kill later fed its partner in
20\% of hunts in \arm{named}, about 10\% in \arm{parent\_id} and \arm{world\_name}, 3\% in
\arm{stable\_id} and 1\% in \arm{unnamed}, so the counts reflect more sharing, not just more
kills.

% Study 4: the main measures in all eleven arms, from figdata/round5.csv.
\begin{figure}[t]
\centering
\begin{tikzpicture}
\begin{groupplot}[
  group style={group size=4 by 1, horizontal sep=0.35cm, yticklabels at=edge left},
  width=0.275\linewidth, height=7.2cm,
  xbar, /pgf/bar width=6pt, xmin=0,
  y dir=reverse, ytick=data, yticklabels from table={figdata/round5.csv}{label},
  yticklabel style={font=\footnotesize},
  xticklabel style={font=\footnotesize, /pgf/number format/fixed},
  scaled x ticks=false,
  title style={font=\small, align=center},
  enlarge y limits=0.06,
  error bars/x dir=both, error bars/x explicit,
  cycle list={{fill=black!30, draw=black!80}},
]
\nextgroupplot[title={Pair hunts\\per run}]
\addplot table[y=x, x=pair_hunts, x error=pair_hunts_se] {figdata/round5.csv};
\nextgroupplot[title={Couples per\\generation}]
\addplot table[y=x, x=couples, x error=couples_se] {figdata/round5.csv};
\nextgroupplot[title={Share given\\to a newborn}]
\addplot table[y=x, x=endowment, x error=endowment_se] {figdata/round5.csv};
\nextgroupplot[title={Share of kills\\shared}]
\addplot table[y=x, x=kill_share, x error=kill_share_se] {figdata/round5.csv};
\end{groupplot}
\end{tikzpicture}
\caption{Study 4, eleven arms on the same 30 fresh seeds: means over seeds with standard errors.
``Share of kills shared'' is the share of successful pair hunts after which the hunter who landed
the kill fed its partner, over runs with at least one successful pair hunt.}
\label{fig:round5}
\end{figure}

\begin{figure}[t]
\centering
\begin{tikzpicture}
\begin{axis}[
  width=0.62\linewidth, height=5.2cm,
  xlabel={Generation}, ylabel={Share of holdings put in},
  xmin=0.5, xmax=8.5, ymin=0, ymax=0.65, xtick={1,...,8},
  ticklabel style={font=\footnotesize}, label style={font=\small},
  legend style={font=\footnotesize, draw=none, fill=none, at={(0.98,0.03)}, anchor=south east},
]
\fill[black!10] (axis cs:0.5,0.4) rectangle (axis cs:8.5,0.6);
\node[font=\footnotesize, anchor=west] at (axis cs:0.6,0.5) {typical human first-round range};
\addplot[thick, mark=*] table[x=generation, y=all] {figdata/commons_by_generation.csv};
\end{axis}
\end{tikzpicture}
\caption{Study 4: the share of holdings put into the commons in each generation, averaged over
all 330 runs. The shaded band is
the 40--60\% of the endowment that human participants typically contribute in the first round of
a standard public-goods game \citep{ledyard1995public,chaudhuri2011sustaining}.}
\label{fig:commons}
\end{figure}


\paragraph{Either half of the parental framing is enough.} Study~3 found that parent-picked IDs
raised couple formation, but in that arm parents both performed the act and every agent was told
that parents had given the IDs. Study~4 separated the two. Parents who picked IDs under a prompt
that never said parents gave them (B1), and agents told that their world-assigned IDs came from
their parents (B2), each formed about as many couples as the full parental arm (2.02 and 2.08,
against 2.22, and 1.40 for plain world-assigned IDs). In the registered reading, either suffices.

\paragraph{Namesaking did not raise investment.} Within each couple in \arm{world\_namesake}, the
parent a child was named after gave no more at birth than its co-parent (C1: 15 seeds each way).
Two weaker leads, a higher endowment than in the control arm (C2) and more later gifts from the
namesake parent (C3), did not survive correction.

\paragraph{Group labels produce in-group choices, and add to individual handles.} Agents paired
(D1) and hunted (D3) within their own arbitrary band more than chance, the minimal-group pattern.
Gifts were too rare in the band arms for D2 to decide. The factionalism hypothesis predicted that a
band would reduce coordination; instead, adding a band to a world-assigned ID doubled pair hunting
(D4: 16.7 vs 8.4), with hunts concentrated inside bands. Without an individual handle, a band did
nothing for coordination (D5: 3.1 vs 3.5 without any label).

\paragraph{The commons sat at a floor.} Agents put 2--6\% of their holdings into the pot in every
arm, and 88--95\% of contribution decisions were zero. No arm differed (E1--E3). As
Figure~\ref{fig:commons} shows, contributions rose slightly over the run, from 1\% in
generation~1 to 5\% in generation~8.

\paragraph{Exploratory findings.} The following are unadjusted and should be read as leads.
\begin{itemize}
  \item \textbf{One sentence about namesakes halved couple formation.} \arm{world\_namesake}
        formed 1.00 couples per generation, against 2.15 in \arm{world\_name\_shown} (lower on
        29 of 30 seeds, $p\le10^{-5}$) and 1.77 without any handle. The drop was already present
        in generation~1, before any child had been born or named: 5.0 of 10 agents asked for a
        partner, against 6.6 ($p=0.0006$). The only difference at that point is one sentence of
        the system prompt saying that the world names every newborn after one of its parents.
        It is not avoidance of partners who share a name: agents listed partners sharing their
        name's base as often as chance predicts ($p=0.41$). The intents of agents who opted out
        read the same in both arms (keep holdings, do not add a rival at the cut), and none
        mentions names.
  \item \textbf{World-assigned IDs lower couple formation below no handle at all.} Stable-ID
        populations formed fewer couples than unnamed ones (1.40 vs 1.77, $p=0.006$), as in
        Study~2 (1.44 vs 1.74, $p=0.046$). Parent-picked IDs formed more (2.22, $p=0.005$).
  \item \textbf{Study~1's named-versus-unnamed couples effect was smaller here} (2.00 vs 1.77,
        $p=0.14$). The contrast that held in every study is a parent-given handle against a
        world-assigned ID.
  \item \textbf{Bands reached into the family.} In \arm{band\_only}, a child's band was drawn by
        lot, independent of its parents. Parents gave a child drawn into the other band less
        than one drawn into their own (0.31 vs 0.36 of holdings; own band higher in 21 of 30
        seeds; $p=0.006$). Both band arms gave less to newborns than their unbanded counterparts
        (0.34 vs 0.40, $p=0.009$; 0.36 vs 0.42, $p=0.027$). Bands appeared in 14\% of
        band-only pairing intents, but almost never in public messages.
  \item \textbf{Kin remained the main target of gifts without names.} Unnamed agents directed
        90\% of their gift amounts to kin, against 52\% in \arm{named} ($p=0.004$), as in Study~1.
\end{itemize}


\section{Comparison with human research}
\label{sec:human}

Study~4 registered each family with the human finding it would be compared with.
Table~\ref{tab:human} sets the results beside those findings. The comparison is of direction only:
the agents face no real stakes, every agent is one model, and none of the quantities is calibrated
to a human measure.

\begin{table}[t]
\centering
\small
\caption{Results set against the human findings they were registered to be compared with.}
\label{tab:human}
\begin{tabularx}{\linewidth}{@{}L L L >{\raggedright\arraybackslash}p{2.1cm}@{}}
\toprule
Human finding & Our analogue & Result & Same direction? \\
\midrule
Identifying individuals raises cooperation \citep{andreoni2004public,christens2019identification} &
Any stable individual handle vs none: pair hunts (A1, A2; Studies 1--2) &
Pair hunting rose two- to fivefold with any handle & Yes \\
\addlinespace
Identifying only groups does little \citep{christens2019identification} &
A band without an individual handle (D5, E2) &
No change in pair hunting or the commons & Yes \\
\addlinespace
Names alone have no direct effect on public-goods contributions \citep{sampaio2023identity} &
Named vs stable ID: the commons (E3) &
No difference, at a floor & Consistent, not informative \\
\addlinespace
Arbitrary labels produce in-group favouritism \citep{tajfel1971social} &
In-band couple, hunt and gift bias (D1--D3); endowment to an other-band child (exploratory) &
In-band couples and hunts; less given to an other-band child & Yes \\
\addlinespace
Shared group labels split LLM agents into factions and degrade coordination \citep{delgiudice2026identity} &
A band added to an individual ID (D4) &
Pair hunting doubled, mostly within bands & No, opposite \\
\addlinespace
Children named after a parent receive more of that parent's investment \citep{saffa2021namesake} &
Namesake parent minus co-parent (C1--C3) &
No difference at birth; weak lead in later gifts & No \\
\addlinespace
First-round public-goods contributions of 40--60\%, declining with repetition \citep{ledyard1995public,chaudhuri2011sustaining} &
Commons contribution share by generation &
About 3\%, rising from 1\% to 5\% & No \\
\addlinespace
Help concentrates on kin \citep{hamilton1964genetical} &
Share of gift amounts to kin (Studies 1 and 4, exploratory) &
Most concentrated without names (84--90\%) & Yes \\
\bottomrule
\end{tabularx}
\end{table}

Four of the eight comparisons agree in direction, one is uninformative, and three disagree. The
agreements cluster around identification and labels: individual handles enable cooperation, group
labels alone do not, and arbitrary group labels bias choices toward the in-group. The
disagreements concern investment and contribution: namesaking does not move parental investment,
and the agents free-ride on the commons almost completely. The commons result carries a caveat
specific to this world. Survival is decided by rank, so anything put into the pot benefits rivals
as much as the contributor, and contributing always lowers an agent's standing relative to those
who hold back. A standard public-goods game has no such tournament.


\section{Discussion}
\label{sec:discussion}

\paragraph{Names do separable jobs.} Across four pre-registered studies, what a handle lets agents
do separated cleanly from what it says about where the agent came from. Any stable individual
handle, whether a parent-chosen name, a world-assigned name, a parent-picked ID or a bare
world-assigned ID, raised pair hunting severalfold over no handle at all, in two independent seed
sets. Pair hunting is the one activity that needs a specific partner on a specific round, and a
handle is what lets agents arrange that across decisions in which their bracket numbers are
reshuffled. This fits the argument that reputation and repeated cooperation presuppose
persistent identity \citep{hu2026dissociative}.

\paragraph{Family formation follows a parental frame, not the label.} Couple formation responded
to whether parents were said to give the handle, not to what the handle was. A meaningless ID
raised couple formation when parents picked it (Studies 3 and 4), when parents picked it but no
prompt said so (B1), and when no parent picked it but agents were told parents had (B2). Either
half of the framing was enough, and the full frame raised couple formation only a little further
(2.22, against 2.02 and 2.08; not tested). Against no handle at all, the
pattern in Study~4 was asymmetric (exploratory): parent-given IDs raised couple formation, while
world-assigned IDs lowered it, as they had in Study~2. Study~1's named-versus-unnamed couples
effect did not reach significance on the fresh seeds, so the contrast to rely on is a
parent-given handle against a world-assigned one. One reading is that a world-assigned ID
supplies an administrative frame in which agents act as individual optimisers, and a parental
frame supplies a family one. The experiment shows that the frame changes behaviour; it does not
show how the model represents it.

\paragraph{Investment in newborns is narrower.} Only one arm consistently raised the share of
holdings parents gave a newborn: names composed by the parents themselves (Study~3 C4, Study~4
A4). Neither half of the parental frame raised it on its own with IDs (0.42 and 0.43, against
0.42 for world-assigned IDs; not tested), and neither did a world-chosen name shown at birth or a
namesake. In the named arm, parents propose a name in the
same decision in which they set the endowment, so composing a name for the child and investing
in it cannot be separated here. That the effect needs this specific act is consistent with
naming as a costly signal of investment, but the design cannot distinguish that from simple
priming within one decision.

\paragraph{A single sentence can move reproduction.} The largest exploratory effect in Study~4
was not one we predicted. Telling agents that the world names each newborn after one of its
parents halved couple formation, from generation~1 onward, before any child existed. The agents'
stated reasons for opting out did not change and never mentioned names. Whatever the mechanism,
two practical points follow. Small, plausible-sounding changes to the description of a world can
have effects as large as the manipulations a study is designed around. And the reasons agents give
for their decisions are not a reliable guide to what moved them.

\paragraph{Group labels: in-group bias without factional breakdown.} Arbitrary Blue and Green
bands reproduced the minimal-group pattern: agents paired and hunted within their band more than
chance, and parents gave less to a child drawn into the other band, so the bias reached into the
family itself. Yet the band did not degrade coordination as
shared labels did in \citet{delgiudice2026identity}. Added to an individual ID, it doubled pair
hunting. A plausible reason is that a pair hunt can succeed with any partner, so a band works as
a focal set that narrows the search for one, while in the tasks where labels caused factionalism
the whole group had to reach a decision together. Without an individual handle, a band did nothing for
coordination, which is the LLM counterpart of the human finding that identifying groups, rather
than individuals, does little for cooperation \citep{christens2019identification}.

\paragraph{The commons.} The agents behaved as near-perfect free-riders, contributing about a
tenth of what people typically contribute and slightly more, not less, as the run went on. Under
rank-based survival, contributing is dominated in relative terms, and Haiku~5.5 at low effort
appears to act on that consistently. With contributions this close to zero, the registered
identifiability tests (E1--E3) had no room to detect anything. Whether the floor comes from the
tournament or from the model would need a commons outside the cut, or a different model.

\paragraph{Implications for multi-agent systems.} Identity scaffolding is not neutral. Whether
agents carry names, who is said to have given them, and whether agents also carry a group label
changed coordination, family formation, investment, sharing and in-group choice by large factors,
holding everything else fixed. Studies of LLM agent societies should treat naming and its
description as a variable to be reported, and ideally controlled.


\section{Limitations}
\label{sec:limits}

\begin{itemize}
  \item \textbf{One model, one framing, one world.} Every agent is Claude Haiku~5.5 at low effort,
        and every arm shares one system prompt apart from its handle clauses. The results are
        claims about that model under that framing. Whether they hold for other models, larger
        efforts or other wordings is untested.
  \item \textbf{The comparisons with people are of direction only.} The simulation has no stakes
        for the agents, and its quantities (couples, endowments, contribution shares) are not
        calibrated to any human measure.
  \item \textbf{Prompt asymmetry.} ID arms are told that an ID is kept for life and never shared;
        name arms are not told this about names. The word-versus-ID contrasts carry that
        sentence.
  \item \textbf{Timing in Studies 2 and 3.} Their runs happened hours after the Study~1 runs they
        are paired with. Study~4 runs every arm of a seed in one window.
  \item \textbf{Correlated measures.} Couples, endowments, newborn survival and displacement are
        measured on the same populations and partly drive one another.
  \item \textbf{The commons is a tournament.} Survival is decided by rank, so contributing to the
        pot always lowers an agent's standing relative to agents who hold back. The near-zero
        contributions may reflect that, not a general reluctance to contribute.
  \item \textbf{Unregistered effects.} The namesake effect on couple formation, the band effect
        on endowments and the comparison of world-assigned IDs with no handle are exploratory.
        They need their own registered test before being relied on.
  \item \textbf{Rare behaviours are underpowered.} Gifts were rare in the band arms, so the
        registered in-band gift test (D2) rested on six seeds. Kill sharing had many seeds with no
        shared kill in either arm.
  \item \textbf{Shared history in every arm.} Because every agent sees its full history with
        every other agent, the experiment cannot show whether handles make reputation possible,
        only whether they change behaviour when information is equal.
\end{itemize}

\section{Conclusion}

In generational populations of LLM agents, a name does at least two separable things. As a
stable handle, it lets agents find the same partner again, and any stable handle does that job
equally well. As something parents are said to give, it frames the population as a family:
either the parental act or merely the wording raised couple formation, while, in exploratory
analyses, a world-assigned ID lowered it below no handle at all. Only names that parents composed themselves raised investment in
newborns. Arbitrary group labels reproduced the minimal-group pattern, yet strengthened rather than
weakened coordination when added to individual handles. Two registered human parallels failed:
namesaking did not raise the namesake parent's investment, and the agents almost entirely
free-rode on a shared pot. One unregistered sentence about namesakes halved couple formation.
Together these results suggest that how agents are named, and how the naming is described, should
be treated as an experimental variable in multi-agent LLM research rather than as scaffolding.


\section*{Acknowledgments}

The simulator, the analysis code, the pre-registration drafts and a first draft of this paper
were written with Claude Opus~5.5 (Anthropic), used as a coding and writing assistant under the
author's direction. The author set the research questions and the original design, approved each
pre-registration before its data were collected, and reviewed and is responsible for the final
text. The agents in the simulation are Claude Haiku~5.5. All API costs were paid by the author.

\section*{Data and code availability}

The simulator, configurations, pre-registrations, event logs of every run and the analysis
outputs are at \url{https://github.com/Vaughanwj/Agent-Names-Test}\ifx\SOFTWAREDOI\empty\else, and
the version this paper reports (v1.0) is archived at Zenodo,
\href{https://doi.org/\SOFTWAREDOI}{doi:\SOFTWAREDOI}\fi. Each pre-registration is
identified by the commit at which it was registered (Table~\ref{tab:studies}). The figures are
generated from the event logs by \texttt{paper/make\_figures.py}. Study~4's registered analysis
is \texttt{scripts/analyze\_round5.py}, and its exploratory analyses are
\texttt{scripts/explore\_round5.py}. The code is released under the MIT licence and the data under
CC BY 4.0.

\section*{License}

\doclicenseThis

\bibliographystyle{plainnat}
\bibliography{refs}

\end{document}
